% !LW recipe = PDFLaTeX Auto
\documentclass[10pt]{article}
\usepackage[top=1.1in, bottom=1.1in, left=1.1in, right=1.1in, letterpaper]{geometry}
\usepackage[numbers,compress]{natbib}
\usepackage{xcolor}
\usepackage{graphicx} % 图片加载
\usepackage{amsmath} %更多数学符号
\usepackage{lmodern}
\usepackage{makecell}  
\usepackage{array} 
\usepackage{tabularx}
% Math packages

\usepackage{mathrsfs}
\usepackage{amsfonts}
\usepackage{amsthm}
\usepackage{bm}
\usepackage{bbm}
\usepackage{esint}
\usepackage{algorithm,algpseudocode}
\usepackage{enumitem} 
\usepackage{pifont}
% Other utility packages
\usepackage{xcolor}
\usepackage{algorithm, algpseudocode}
\usepackage{enumitem}
% \usepackage[utf8]{inputenc}
\usepackage{url}
\usepackage{nicefrac} % compact symbols for 1/2, etc.
\usepackage{textcomp}
\usepackage{setspace}
\usepackage{latexsym}
\usepackage{alltt}
\usepackage{xspace}
\usepackage{subcaption}
\usepackage{booktabs}
\usepackage{multirow} 
\usepackage{float}
\usepackage{appendix}  % 专门处理附录的包
\usepackage{chngcntr} 
\usepackage{threeparttable}
\usepackage{enumitem}
% Hyperref setup
% 定理环境设置
% \newtheorem{theorem}{Theorem}[section]
% \newtheorem{lemma}[theorem]{Lemma}
% \newtheorem{corollary}[theorem]{Corollary}
% \newtheorem{proposition}[theorem]{Proposition}

% \theoremstyle{definition}
% \newtheorem{definition}[theorem]{Definition}
% \newtheorem{example}[theorem]{Example}
% \theoremstyle{remark}
% \newtheorem{remark}[theorem]{Remark}
% \newtheorem{theorem}{Theorem}[section]
% \newtheorem{lemma}{Lemma}[section]
% \newtheorem{corollary}{Corollary}[section]
% \newtheorem{proposition}{Proposition}[section]
% \theoremstyle{definition}
% \newtheorem{definition}{Definition}[section]
% \newtheorem{example}{Example}[section]
% \theoremstyle{remark}
% \newtheorem{remark}{Remark}[section]

\newcommand{\theHalgorithm}{\arabic{algorithm}}
\usepackage{cleveref}
\usepackage{subcaption}
%-----------------------------新命令-------------------------------
\newcommand{\eucalB}{\EuScript{B}}
\newcommand{\eucalP}{\EuScript{P}}
\newcommand{\Var}{\mathrm{Var}}
\newcommand{\Cov}{\mathrm{Cov}}
\newcommand{\Ex}{\mathsf{E}}
\newcommand{\Prob}{\mathsf{P}}
\newcommand{\one}{\mathbbm{1}}
\newcommand{\One}{\bm{1}}
\newcommand{\zero}{\bm{0}}
\newcommand{\RR}{\mathbb{R}}
\newcommand{\ZZ}{\mathbb{Z}}
\newcommand{\EE}{\mathbb{E}}
\newcommand{\PP}{\mathbb{P}}
\newcommand{\Rp}{\RR_+}
\newcommand{\TT}{\mathbb{T}}
\newcommand{\NN}{\mathbb{N}}
\newcommand{\VV}{\mathbb{V}}
\newcommand{\diff}{\mathrm{d}}
\newcommand{\prox}{\mathrm{prox}}
\newcommand{\rprox}{\mathrm{rprox}}
\newcommand{\ran}{\mathrm{ran}}
\newcommand{\Id}{\mathrm{Id}}
\newcommand{\e}{\mathrm{e}}
\newcommand{\bA}{\bm{A}}
\newcommand{\bB}{\bm{B}}
\newcommand{\bH}{\bm{H}}
\newcommand{\bF}{\bm{F}}
\newcommand{\bX}{\bm{X}}
\newcommand{\bx}{\bm{x}}
\newcommand{\bb}{\bm{b}}
\newcommand{\bc}{\bm{c}}
\newcommand{\br}{\bm{r}}
\newcommand{\be}{\bm{e}}
\newcommand{\bbeta}{\bm{\rho}}
\newcommand{\bt}{\bm{t}}
\newcommand{\bh}{\bm{h}}
\newcommand{\bu}{\bm{u}}
\newcommand{\bz}{\bm{z}}
\newcommand{\ba}{\bm{a}}
\newcommand{\by}{\bm{y}}
\newcommand{\bU}{\bm{U}}
\newcommand{\bZ}{\bm{Z}}
\newcommand{\bY}{\bm{Y}}
\newcommand{\blmbd}{\bm{\lambda}}
\newcommand{\bLmbd}{\bm{\Lambda}}
\newcommand{\bg}{\bm{g}}
\newcommand{\bd}{\bm{d}}
\newcommand{\bxi}{\bm{\xi}}
\newcommand{\bdelta}{\bm{\delta}}
\DeclareMathOperator*{\argmin}{\mathrm{argmin}}
\DeclareMathOperator*{\Argmin}{\mathrm{Argmin}}
\DeclareMathOperator*{\minimize}{\mathrm{minimize}\,}
\DeclareMathOperator*{\maximize}{\mathrm{maximize}}
\DeclareMathOperator*{\argmax}{\mathrm{argmax}}
\DeclareMathOperator*{\Diag}{\mathrm{Diag}}
\newcommand{\calX}{\mathcal{X}}
\newcommand{\calY}{\mathcal{Y}}
\newcommand{\calB}{\mathcal{B}}
\newcommand{\calL}{\mathcal{L}}
\newcommand{\calG}{\mathcal{G}}
\newcommand{\calS}{\mathcal{S}}
\newcommand{\calC}{\mathcal{C}}
\newcommand{\calJ}{\mathcal{J}}
\newcommand{\calF}{\mathcal{F}}
\newcommand{\calM}{\mathcal{M}}
\newcommand{\calO}{\mathcal{O}}
\newcommand{\calI}{\mathcal{I}}
\newcommand{\calE}{\mathcal{E}}
\newcommand{\calH}{\mathcal{H}}
\newcommand{\ent}{\mathrm{ent}}
\newcommand{\diag}{\mathrm{diag}}
\newcommand{\dom}{\mathrm{dom}}
\newcommand{\interior}{\mathrm{int}}
\newcommand{\zer}{\mathrm{zer}}
\newcommand{\epi}{\mathrm{epi}}
\newcommand{\amax}{a_{\max}}
\newcommand{\amin}{a_{\min}}
\newcommand{\bmax}{b_{\max}}
\newcommand{\bmin}{b_{\min}}
\newcommand{\bp}{\bm{p}}
\newcommand{\bmu}{\bm{\mu}}
\newcommand{\bD}{\bm{D}}
\newcommand{\mW}{\mathrm{W}}
\newcommand{\bv}{\bm{v}}
\newcommand{\bw}{\bm{w}}
\newcommand{\ta}{\tilde{a}}
\newcommand{\tb}{\tilde{b}}
\newcommand{\tp}{\widetilde{\bp}}
\newcommand{\sumn}{\sum_{i=1}^n}
\newcommand{\sumQ}{\sum_{q=1}^Q}
\newcommand{\sumq}{\sum_{k=1}^q}
\newcommand{\summ}{\sum_{i=1}^m}
\newcommand{\sumd}{\sum_{i=1}^d}
\newcommand{\sumT}{\sum_{k=0}^{T-1}}
\newcommand{\sumK}{\sum_{k=0}^{K-1}}
\newcommand{\sumN}{\sum_{k=1}^{N}}
\newcommand{\bC}{\bm{C}}
\newcommand{\bI}{\bm{I}}
\newcommand{\bM}{\bm{M}}
\newcommand{\bR}{\bm{R}}
\newcommand{\bS}{\bm{S}}
\newcommand{\bV}{\bm{V}}
\newcommand{\bW}{\bm{W}}
\newcommand{\blambda}{\bm{\lambda}}
\newcommand{\tmu}{\widetilde{\bmu}}
\newcommand{\tU}{\widetilde{\bU}}
\newcommand{\tV}{\widetilde{\bV}}
\newcommand{\tM}{\widetilde{\bM}}
\newcommand{\bDelta}{\bm{\Delta}}
\newcommand{\bs}{\bm{s}}
\newcommand{\bzeta}{\bm{\zeta}}
\newcommand{\ox}{\overline{x}}
\newcommand{\ob}{\overline{b}}
\newcommand{\op}{\overline{p}}
\newcommand{\obp}{\overline{\bp}}
\newcommand{\obmu}{\overline{\bmu}}
\newcommand{\obC}{\overline{\bC}}
\newcommand{\obD}{\overline{\bD}}
\newcommand{\bsigma}{\bm{\sigma}}
\newcommand{\bomega}{\bm{\omega}}
\newcommand{\tomega}{\widetilde{\omega}}
\newcommand{\tr}{\mathrm{tr}}
\newcommand{\an}{\mathrm{an}}
\newcommand{\bj}{\bm{j}}
\newcommand{\bG}{\bm{G}}
\newcommand{\bcalG}{\bm{\calG}}
\newcommand{\bcalK}{\bm{\calK}}
\newcommand{\bcalX}{\bm{\calX}}
\newcommand{\fX}{\mathfrak{X}}
\newcommand{\bfX}{\bm{\fX}}
\newcommand{\bell}{\bm{\ell}}
\newcommand{\bL}{\bm{L}}
\renewcommand{\vec}{\mathrm{vec}}
\newcommand{\hy}{\widehat{y}}
\newcommand{\ul}[1]{\underline{#1}}
\newcommand{\xmark}{\ding{55}}
\newtheorem{assumption}{Assumption}
\newtheorem{Def}{Definition}
% \newtheorem{corollary}{Corollary}
\newenvironment{proofsketch}
    {\noindent\textit{Proof sketch.}}
    {\hfill$\square$}
% \newtheorem{Proposition}{proposition} 
% \crefname{Proposition}{proposition}{propositions}
% \Crefname{Proposition}{Proposition}{Propositions}
\newcommand{\eproof}{\hfill $\square$}
\newcommand{\bcL}{\overline{\calL}}
\newcommand{\KL}{K\L\xspace}
\newcommand{\res}{\mathrm{res}}
\newtheorem{lemma}{Lemma}

\newcommand{\commyy}[1]{{\color{red}(yy: #1)}} % Yuan Yao's comments
\newcommand{\blue}[1]{\begin{color}{blue}#1\end{color}} 
\newcommand{\red}[1]{\begin{color}{red}#1\end{color}}

\numberwithin{equation}{section}
\Crefname{algorithm}{Algorithm}{Algorithms}



\title{Dual-Smoothing Augmented Lagrangian Methods for Deterministic Nonconvex Equality-Constrained Optimization}
\author{}

\begin{document}

\maketitle

\section{Introduction}
\label{sec:introduction}

We study first-order augmented-Lagrangian-type methods for smooth nonconvex
optimization with nonlinear equality constraints.  For this problem class,
an $O(\epsilon^{-3})$ first-order complexity is a standard benchmark for
penalty and augmented-Lagrangian frameworks that do not rely on solving
nonconvex augmented-Lagrangian subproblems to high accuracy.  Recent
linearized quadratic penalty methods can improve this benchmark under a
different Gauss--Newton-type model, and restoration or SQP-type frameworks
may obtain faster rates under different algorithmic structures.  Our aim is
to develop a single-loop dual-smoothing ALM framework whose iterations are
explicit first-order steps, and to establish a weighted finite-length
estimate that yields the standard $O(\epsilon^{-3})$ benchmark under compact
working-set and uniform LICQ assumptions.

The main technical idea is to introduce a dual proximal smoothing term in
the multiplier variable and then eliminate the smoothed multiplier in closed
form.  The anchor multiplier $\bmu^k$ is the evolving multiplier, whereas the
virtual multiplier $\blmbd^k$ is the maximizer induced by the dual-smoothing
term.  On the algorithmic trajectory, this yields the reduced potential
\[
    \Phi_k(\bx;\bmu)
    =
    f(\bx)
    +
    \langle \bmu,\bc(\bx)\rangle
    +
    \frac{\bar\rho_k}{2}\|\bc(\bx)\|^2,
    \qquad
    \bar\rho_k=\rho+s_k .
\]
Thus the primal step is a gradient step on a moving reduced augmented
Lagrangian with anchor $\bmu^k$.  This viewpoint separates the stability of
the multiplier anchor from the effective penalty strength.

The intended contributions are as follows.
\begin{itemize}
    \item We formulate a dual-smoothing ALM from a proximal-perturbed
    Lagrangian whose smoothed multiplier has a closed-form relation.
    \item Under a compact working-set assumption and a uniform LICQ
    condition, we derive a reduced-potential descent mechanism in which the
    smoothness depends on the bounded auxiliary multiplier $\bmu^k$, rather
    than requiring a priori boundedness of the virtual multiplier
    $\blmbd^k$.
    \item We use $\bar\rho_k=\rho+s_k$ as the scheduled effective penalty
    parameter induced by dual smoothing and outline a residual-adaptive
    damped anchor update.  Under the scaling
    $\bar\rho_k\asymp k^{1/3}$ and $\eta_k\asymp k^{-1/3}$, we prove a
    weighted finite-length estimate that yields the balanced
    $O(\epsilon^{-3})$ complexity.
    \item The approximate KKT residual is naturally stated with the effective
    multiplier
    $\by^k:=\bmu^k+\bar\rho_k\bc(\bx^k)$, while the virtual
    dual-smoothing multiplier satisfies
    $\by^k=\blmbd^k+\rho\bc(\bx^k)$, so their difference is controlled
    directly by the feasibility residual.
\end{itemize}

\section{Problem Setup and Algorithm}
\label{sec:method}

We consider the following deterministic nonconvex optimization problem over
the full Euclidean space with smooth objective and nonlinear equality
constraints:
\begin{align}
    \min_{\bx\in\RR^n}\quad & f(\bx) \notag\\
    \text{s.t.}\quad & \bc(\bx)=\zero, \tag{\textbf{P}}
\end{align}
where $f:\RR^n\to\RR$ is continuously differentiable, and
$\bc:\RR^n\to\RR^m$ is continuously differentiable.

Fix a base penalty $\rho>0$ and a dual-smoothing scale $s_k>0$.  At
iteration $k$, define the proximal-perturbed dual-smoothing potential
\begin{equation}\label{eq:dual-smoothing-pplf}
    \mathcal P_k(\bx,\blmbd;\bmu)
    :=
    f(\bx)
    +\langle \blmbd,\bc(\bx)\rangle
    +\frac{\rho}{2}\|\bc(\bx)\|^2
    -\frac{1}{2s_k}\|\blmbd-\bmu\|^2 .
\end{equation}
The map $\blmbd\mapsto\mathcal P_k(\bx,\blmbd;\bmu)$ is strongly concave.
Its first-order condition gives the unique smoothed multiplier
\begin{equation}\label{eq:lambda-star-sk}
    \blmbd_k(\bx,\bmu)
    =
    \bmu+s_k\bc(\bx).
\end{equation}
Substituting \eqref{eq:lambda-star-sk} into
\eqref{eq:dual-smoothing-pplf} gives the reduced augmented Lagrangian
\begin{equation*}
    \Phi_k(\bx;\bmu)
    :=
    f(\bx)
    +
    \langle \bmu,\bc(\bx)\rangle
    +
    \frac{\bar\rho_k}{2}\|\bc(\bx)\|^2,
    \qquad
    \bar\rho_k:=\rho+s_k.
\end{equation*}
The effective KKT multiplier associated with the reduced potential at
$(\bx^k,\bmu^k)$ is
\begin{equation}\label{eq:effective-multiplier}
    \by^k
    :=
    \bmu^k+\bar\rho_k\bc(\bx^k).
\end{equation}

\paragraph{Algorithmic updates.}
Given $(\bx^k,\bmu^k,\by^k)$ with $\by^k$ defined by
\eqref{eq:effective-multiplier}, the proposed deterministic adaptive
dual-smoothing ALM performs
\begin{align}
    \bx^{k+1}
    &\in
    \argmin_{\bx\in\RR^n}
    \left\{
    \left\langle
    \nabla f(\bx^k)
    +J_{\bc}(\bx^k)^\top\by^k,
    \bx-\bx^k
    \right\rangle
    +\frac{1}{2\eta_k}\|\bx-\bx^k\|^2
    \right\},
    \label{eq:x-update-det}\\
    \bmu^{k+1}
    &=
    \bmu^k+\tau_k\bc(\bx^k),
    \label{eq:mu-update-det}\\
    \by^{k+1}
    &=
    \bmu^{k+1}+\bar\rho_{k+1}\bc(\bx^{k+1}),
    \label{eq:y-update-det}
\end{align}
where $\eta_k>0$ is the primal stepsize and $\tau_k\ge0$ is the effective
dual stepsize.  The method only requires $s_k>0$ and
$0\le\tau_k\le s_k$ for the interpolation interpretation below; concrete
parameter schedules are left unspecified at the algorithmic level.

\begin{algorithm}[H]
    \caption{Adaptive Dual-Smoothing ALM for Smooth Equality Constraints}
    \label{alg:adaptive-dsalm}
    \small
    \begin{algorithmic}[1]
        \STATE \textbf{Input:} $\bx^0\in\RR^n$, $\bmu^0\in\RR^m$,
        base penalty $\rho>0$, smoothing scales $\{s_k\}_{k\ge0}$,
        primal stepsizes $\{\eta_k\}_{k\ge0}$, and effective dual stepsizes
        $\{\tau_k\}_{k\ge0}$.
        \STATE Set $\bar\rho_0=\rho+s_0$ and define $\by^0$ by
        \eqref{eq:effective-multiplier}.
        \FOR{$k=0,1,\ldots,K-1$}
            \STATE Update $\bx^{k+1}$ by \eqref{eq:x-update-det}.
            \STATE Update $\bmu^{k+1}$ by \eqref{eq:mu-update-det}.
            \STATE Set $\bar\rho_{k+1}=\rho+s_{k+1}$.
            \STATE Update $\by^{k+1}$ by \eqref{eq:y-update-det}.
        \ENDFOR
        \STATE \textbf{Output:} for $K\ge2$, choose
        $\kappa\in
        \argmin_{k\in\{\lceil K/2\rceil,\ldots,K-1\}}
        \bar\rho_k\|\Delta\bx^{k+1}\|^2$ and return
        $(\bx^\kappa,\by^\kappa)$.
    \end{algorithmic}
\end{algorithm}

\begin{remark}[Equivalent explicit primal step]
\label{rmk:explicit-primal-step}
The primal update \eqref{eq:x-update-det} reduces to the explicit gradient
step
\[
    \bx^{k+1}
    =
    \bx^k
    -
    \eta_k
    \Bigl(
    \nabla f(\bx^k)
    +J_{\bc}(\bx^k)^\top\by^k
    \Bigr).
\]
Thus the method is fully explicit and single-loop in the primal variable.
\end{remark}

\begin{remark}[Dual-smoothing interpolation]
\label{rmk:dual-smoothing-interpolation}
The virtual smoothed multiplier at iteration $k$ is
\[
    \blmbd^k
    :=
    \bmu^k+s_k\bc(\bx^k).
\]
Since $\tau_k\le s_k$, defining
\[
    \theta_k:=\frac{\tau_k}{s_k}\in[0,1]
\]
gives
\[
    \bmu^{k+1}
    =
    (1-\theta_k)\bmu^k+\theta_k\blmbd^k.
\]
Thus the anchor update is a convex interpolation between the old anchor and
the smoothed multiplier, while the analysis uses the effective dual stepsize
$\tau_k$.
\end{remark}

\begin{remark}[Reduced potential viewpoint]
\label{rmk:fixed-penalty-view}
By \eqref{eq:lambda-star-sk}, the virtual multiplier satisfies
\[
    \blmbd^k=\bmu^k+s_k\bc(\bx^k).
\]
Combining this identity with \eqref{eq:effective-multiplier} yields
\[
    \by^k
    =
    \blmbd^k+\rho\,\bc(\bx^k)
    =
    \bmu^k+\bar\rho_k\bc(\bx^k).
\]
Therefore the primal direction in \eqref{eq:x-update-det} is exactly
\[
    \nabla_{\bx}\Phi_k(\bx^k;\bmu^k)
    =
    \nabla f(\bx^k)+J_{\bc}(\bx^k)^\top\by^k.
\]
The dual-smoothing term is the source of the effective penalty
$\bar\rho_k=\rho+s_k$; the parameter $\tau_k$ only controls the actual
movement of the anchor multiplier.
\end{remark}

\section{Assumptions}
\label{sec:assumptions}

We next state the assumptions used in the deterministic analysis.  The
analysis is conditional on the generated primal iterates remaining in a
compact working set.  We do not impose any a priori boundedness assumption
on the generated dual sequences $\{\blmbd^k\}$ and $\{\bmu^k\}$, and we do
not use any local or global error bound condition.  Instead, boundedness of
the anchor sequence $\{\bmu^k\}$ will follow from a finite-variation anchor
update.

\begin{assumption}
\label{ass:working-set}
The problem is posed over $\RR^n$.  There exists a compact convex set
$S\subset\RR^n$ such that the iterates generated by
Algorithm~\ref{alg:adaptive-dsalm} satisfy
\[
    \bx^k\in S,
    \qquad \forall k\ge 0.
\]
\end{assumption}

\begin{assumption}
\label{ass:smoothness}
The function $f$ is $L_f$-smooth on the compact convex set $S$, namely
\[
    \|\nabla f(\bx)-\nabla f(\by)\|
    \le
    L_f\|\bx-\by\|,
    \qquad \forall \bx,\by\in S.
\]
The mapping $\bc$ is continuously differentiable on $\RR^n$, and there exist
constants $C_c,G_c,L_c>0$ such that, for all $\bx,\by\in S$,
\[
    \|\bc(\bx)\|\le C_c,
    \qquad
    \|J_{\bc}(\bx)\|\le G_c,
\]
and
\[
    \|J_{\bc}(\bx)-J_{\bc}(\by)\|
    \le
    L_c\|\bx-\by\|.
\]
\end{assumption}

\begin{assumption}
\label{ass:strong-licq}
There exists a constant $\nu>0$ such that
\[
    \sigma_{\min}\bigl(J_{\bc}(\bx)\bigr)\ge \nu,
    \qquad
    \forall \bx\in S.
\]
Equivalently,
\[
    \|J_{\bc}(\bx)^\top \by\|
    \ge
    \nu \|\by\|,
    \qquad \forall \bx\in S,\ \forall \by\in\RR^m.
\]
In particular, $m\le n$ and $J_{\bc}(\bx)$ has full row rank on $S$.
\end{assumption}

Throughout the analysis, define
\[
    M_f:=\sup_{\bx\in S}\|\nabla f(\bx)\|,
    \qquad
    D_S:=\sup_{\bx,\by\in S}\|\bx-\by\|.
\]

\section{Auxiliary Lemmas}
\label{sec:basic-descent}

For convenience, define the successive differences
\[
    \Delta\bx^{k+1}:=\bx^{k+1}-\bx^k,
    \qquad
    \Delta\bmu^{k+1}:=\bmu^{k+1}-\bmu^k.
\]
Throughout the analysis, the virtual multiplier is defined by
\[
    \blmbd^k:=\bmu^k+s_k\bc(\bx^k).
\]

\begin{lemma}
\label{lem:basic-dual-relations}
Let $\{(\bx^k,\bmu^k)\}_{k\ge 0}$ be generated by
Algorithm~\ref{alg:adaptive-dsalm}.
Then, for all $k\ge0$,
\begin{align}
    \blmbd^k-\bmu^k
    &=
    s_k\bc(\bx^k),
    \label{eq:dual-gap-identity}\\
    % \by^k
    % &=
    % \blmbd^k+\rho\,\bc(\bx^k)
    % =
    % \bmu^k+\bar\rho_k\bc(\bx^k),
    % \notag\\
    \Delta\bmu^{k+1}
    &=
    \tau_k\bc(\bx^k)
    =
    \frac{\tau_k}{s_k}(\blmbd^k-\bmu^k).
    \label{eq:mu-drift-tau}
\end{align}
\end{lemma}

\begin{proof}
The first identity is the closed-form maximizer relation
\eqref{eq:lambda-star-sk} evaluated at $(\bx^k,\bmu^k)$.  The second identity
is exactly \eqref{eq:mu-update-det} combined with
\eqref{eq:dual-gap-identity}.
\end{proof}

We first state the following parameter conditions.  Let $B_\tau>0$ be a
fixed damping constant, with the convenient scaling
$B_\tau\asymp\max\{1,\|\bc(\bx^0)\|\}$, and let
$\{\bar\tau_k\}_{k\ge0}$ be auxiliary caps satisfying
$0\le \bar\tau_k\le s_k$.  For every $k\ge0$, choose
\begin{equation}
\label{eq:adaptive-anchor-stepsize}
    \tau_k
    =
    \min\left\{
    \frac{B_\tau}
    {\|\bc(\bx^k)\|(k+1)^2\log(k+2)},
    \bar\tau_k
    \right\},
\end{equation}
where the first term inside the minimum is interpreted as $+\infty$ when
$\bc(\bx^k)=\zero$.  Only $0\le\bar\tau_k\le s_k$ is required here, and its
concrete form will be specified later from the descent fluctuation estimates.

\begin{remark}[Choice of the damping budget]
The constant $B_\tau$ controls the total-variation budget of the anchor
multiplier.  A convenient practical normalization is
\[
    B_\tau
    \asymp
    \max\{1,\|\bc(\bx^0)\|\}.
\]
This choice ties the damping scale to the initial feasibility violation and
keeps a strictly positive update budget when $\bc(\bx^0)=\zero$.  The
convergence proof only requires $B_\tau>0$; this normalization affects only
the constants through $M_\mu$.
\end{remark}

\begin{lemma}
\label{lem:mu-bounded-summable-damping}
Let $\{(\bx^k,\bmu^k)\}_{k\ge 0}$ be generated by
Algorithm~\ref{alg:adaptive-dsalm}.  Assume that the dual stepsizes satisfy
\eqref{eq:adaptive-anchor-stepsize}.
Define
\[
    M_\mu
    :=
    \|\bmu^0\|
    +
    B_\tau
    \sum_{t=0}^{\infty}
    \frac{1}{(t+1)^2\log(t+2)}.
\]
The series in the definition is finite, hence $M_\mu<\infty$.
Then
\[
    \sum_{k=0}^{\infty}\|\bmu^{k+1}-\bmu^k\|
    <
    \infty,
    \qquad
    \|\bmu^k\|\le M_\mu,
    \qquad
    \forall k\ge0.
\]
\end{lemma}

\begin{proof}
This proof only uses the damping bound implied by
\eqref{eq:adaptive-anchor-stepsize}; it does not require the concrete form of
the cap $\bar\tau_k$.  By
\eqref{eq:adaptive-anchor-stepsize}, for every $k\ge0$,
\[
    \tau_k\|\bc(\bx^k)\|
    \le
    \frac{B_\tau}{(k+1)^2\log(k+2)}.
\]
The inequality is immediate when $\bc(\bx^k)\ne\zero$, and it is trivial
when $\bc(\bx^k)=\zero$.  Since
\[
    \sum_{k=0}^{\infty}
    \frac{1}{(k+1)^2\log(k+2)}
    \le
    \frac{1}{\log 2}
    \sum_{k=0}^{\infty}\frac{1}{(k+1)^2}
    <
    \infty,
\]
using \eqref{eq:mu-drift-tau} gives
\[
    \sum_{k=0}^{\infty}\|\bmu^{k+1}-\bmu^k\|
    =
    \sum_{k=0}^{\infty}
    \tau_k\|\bc(\bx^k)\|
    \le
    B_\tau
    \sum_{k=0}^{\infty}
    \frac{1}{(k+1)^2\log(k+2)}
    <
    \infty.
\]
For any $k\ge0$, telescoping yields
\[
    \bmu^k
    =
    \bmu^0
    +
    \sum_{t=0}^{k-1}
    \left(\bmu^{t+1}-\bmu^t\right).
\]
Taking norms and using the preceding summability estimate gives
\[
    \|\bmu^k\|
    \le
    \|\bmu^0\|
    +
    \sum_{t=0}^{k-1}
    \|\bmu^{t+1}-\bmu^t\|
    \le
    M_\mu.
\]
\end{proof}

For the rate statements below, we impose the following stepsize compatibility
condition:
\[
    \eta_k = \frac{1}{C_\eta\bar\rho_k}.
\]
The sufficient descent condition is guaranteed by choosing
\[
    C_\eta
    \ge
    L_cC_c+G_c^2+1
    +
    \frac{L_f+L_cM_\mu}{\bar\rho_0}.
\]
\begin{remark}[Choice of the initial effective penalty]
A convenient parameter choice is to take the base penalty so that
\[
    \bar\rho_0=\rho+s_0
    \gtrsim
    \max\{1,L_f+L_cM_\mu\}.
\]
This keeps the last term in the lower bound for $C_\eta$ uniformly
controlled.  This recommendation is a constant-level normalization: the
$O(\epsilon^{-3})$ iteration-complexity exponent is still determined by the
schedule $s_k=c_s(k+1)^{1/3}$.
\end{remark}

\begin{lemma}
\label{lem:effective-multiplier-feasibility-bound}
Suppose that Assumptions~\ref{ass:working-set}--\ref{ass:strong-licq} hold.
Assume further that the dual stepsizes satisfy
\eqref{eq:adaptive-anchor-stepsize}.
Let $M_\mu$ be defined in
Lemma~\ref{lem:mu-bounded-summable-damping}.  Then, for every $k\ge0$,
\[
    \|\bc(\bx^k)\|
    \le
    \frac{1}{\bar\rho_k}
    \left(
    \frac{M_f}{\nu}+M_\mu
    \right)
    +
    \frac{1}{\eta_k\nu\bar\rho_k}
    \|\Delta\bx^{k+1}\|
\]
\end{lemma}

\begin{proof}
By Lemma~\ref{lem:mu-bounded-summable-damping}, condition
\eqref{eq:adaptive-anchor-stepsize} gives $\|\bmu^k\|\le M_\mu$ for all
$k\ge0$.
Since the problem is posed over $\RR^n$, the primal update
\eqref{eq:x-update-det} gives
\[
    \frac{1}{\eta_k}(\bx^k-\bx^{k+1})
    =
    \nabla f(\bx^k)
    +
    J_{\bc}(\bx^k)^\top\by^k.
\]
Equivalently,
\[
    J_{\bc}(\bx^k)^\top\by^k
    =
    \frac{1}{\eta_k}(\bx^k-\bx^{k+1})
    -
    \nabla f(\bx^k).
\]
Taking norms and using $\|\nabla f(\bx^k)\|\le M_f$ gives
\[
    \|J_{\bc}(\bx^k)^\top\by^k\|
    \le
    \frac{1}{\eta_k}\|\Delta\bx^{k+1}\|
    +
    M_f.
\]
By Assumption~\ref{ass:strong-licq},
\[
    \nu\|\by^k\|
    \le
    \|J_{\bc}(\bx^k)^\top\by^k\|.
\]
Therefore,
\[
    \|\by^k\|
    \le
    \frac{1}{\nu}
    \left(
    \frac{1}{\eta_k}\|\Delta\bx^{k+1}\|+M_f
    \right).
\]
Since $\bar\rho_k\bc(\bx^k)=\by^k-\bmu^k$, we have
\[
    \bar\rho_k\|\bc(\bx^k)\|
    \le
    \|\by^k\|+\|\bmu^k\|.
\]
Using $\|\bmu^k\|\le M_\mu$ and substituting the preceding bound for
$\|\by^k\|$ proves the claim.
\end{proof}

\begin{lemma}
\label{lem:descent-skeleton}
Suppose that the conditions of \Cref{lem:effective-multiplier-feasibility-bound}
hold, and define
\begin{equation*}
    L_k:=L_f+L_cM_\mu+\bar\rho_k\bigl(L_c C_c+G_c^2\bigr) = \calO(\bar\rho_k).
\end{equation*}

If, in addition,
\begin{equation}\label{eq:stepsize-sufficient-descent}
    \frac{1}{2\eta_k}-\frac{L_k}{2}\ge \frac{1}{2}\bar\rho_k,
\end{equation}
then, for every $k\ge0$, it holds that
\begin{equation}\label{eq:variable-penalty-descent-sufficient}
\begin{aligned}
    \Phi_{k+1}(\bx^{k+1};\bmu^{k+1})-\Phi_k(\bx^k;\bmu^k)
    &\le -\frac{1}{2}\bar\rho_k\|\Delta\bx^{k+1}\|^2
    +C_c\|\Delta\bmu^{k+1}\| +\frac{\bar\rho_{k+1}-\bar\rho_k}{2}\|\bc(\bx^{k+1})\|^2 .
\end{aligned}
\end{equation}
\end{lemma}

\begin{proof}
Split the fixed-penalty change as
\[
\begin{aligned}
    \Phi_k(\bx^{k+1};\bmu^{k+1})-\Phi_k(\bx^k;\bmu^k)
    &=
    \Phi_k(\bx^{k+1};\bmu^k)-\Phi_k(\bx^k;\bmu^k) +
    \Phi_k(\bx^{k+1};\bmu^{k+1})-\Phi_k(\bx^{k+1};\bmu^k).
\end{aligned}
\]
Since $S$ is convex, the segment between $\bx^k$ and $\bx^{k+1}$ lies in
$S$.  For fixed $\bmu^k$, the function $\Phi_k(\cdot;\bmu^k)$ is smooth on
$S$ with smoothness constant bounded by $L_k$, and
\[
    \nabla_{\bx}\Phi_k(\bx^k;\bmu^k)
    =
    \nabla f(\bx^k)+J_{\bc}(\bx^k)^\top\by^k .
\]
The optimality of \eqref{eq:x-update-det}, with $\bx=\bx^k$, gives
\[
    \left\langle \nabla_{\bx}\Phi_k(\bx^k;\bmu^k),\Delta\bx^{k+1}\right\rangle
    \le
    -\frac{1}{2\eta_k}\|\Delta\bx^{k+1}\|^2 .
\]
The descent lemma gives the primal descent term.  The anchor-drift term obeys
\[
\begin{aligned}
    \Phi_k(\bx^{k+1};\bmu^{k+1})-\Phi_k(\bx^{k+1};\bmu^k)
    &=
    \langle \Delta\bmu^{k+1},\bc(\bx^{k+1})\rangle\\
    &\le C_c\|\Delta\bmu^{k+1}\|.
\end{aligned}
\]
Combining these two bounds gives the fixed-penalty descent bound.  Finally,
\[
    \Phi_{k+1}(\bx^{k+1};\bmu^{k+1})
    -
    \Phi_k(\bx^{k+1};\bmu^{k+1})
    =
    \frac{\bar\rho_{k+1}-\bar\rho_k}{2}\|\bc(\bx^{k+1})\|^2 ,
\]
because only the penalty coefficient changes from $\Phi_k$ to $\Phi_{k+1}$.
Adding this identity to the fixed-penalty descent bound and using
\eqref{eq:stepsize-sufficient-descent} gives
\eqref{eq:variable-penalty-descent-sufficient}.
\end{proof}

\begin{lemma}
\label{lem:weighted-finite-length}
Suppose that the conditions of Lemma~\ref{lem:descent-skeleton} hold.
Assume in addition that there exists a constant $C_\eta>0$ such that
\[
    \eta_k^{-1}\le C_\eta\bar\rho_k,
    \qquad
    \bar\rho_k=\rho+c_s(k+1)^a,
    \qquad
    c_s>0,\quad a\in(0,1).
\]
Then there exists a constant $C_{\mathrm{desc}}>0$, independent of $K$,
such that, for every $K\ge1$,
\begin{equation}\label{eq:weighted-finite-length}
    \sum_{k=0}^{K-1}
    \bar\rho_k\|\Delta\bx^{k+1}\|^2
    \le
    C_{\mathrm{desc}} .
\end{equation}
\end{lemma}

\begin{proof}
Throughout the proof, $C>0$ denotes a generic constant independent of $k$
and $K$.  The value of $C$ may change from line to line and depends only on
the problem data and fixed algorithmic parameters.
By \eqref{eq:variable-penalty-descent-sufficient},
\[
\begin{aligned}
    \Phi_{k+1}(\bx^{k+1};\bmu^{k+1})-\Phi_k(\bx^k;\bmu^k)
    &\le
    -\frac{1}{2}\bar\rho_k\|\Delta\bx^{k+1}\|^2
    +C_c\|\Delta\bmu^{k+1}\| +
    \frac{\bar\rho_{k+1}-\bar\rho_k}{2}\|\bc(\bx^{k+1})\|^2 .
\end{aligned}
\]
We first control the constraint violation at $\bx^{k+1}$.  By
Lemma~\ref{lem:effective-multiplier-feasibility-bound} and
$\eta_k^{-1}\le C_\eta\bar\rho_k$,
\[
    \|\bc(\bx^k)\|
    \le
    C\left(
    \frac{1}{\bar\rho_k}+\|\Delta\bx^{k+1}\|
    \right).
\]
Since $S$ is convex and $\|J_{\bc}(\bx)\|\le G_c$ on $S$, the mean-value
inequality gives
\[
    \|\bc(\bx^{k+1})-\bc(\bx^k)\|
    \le
    G_c\|\Delta\bx^{k+1}\|.
\]
Therefore
\begin{equation}\label{eq:next-feasibility-from-step}
    \|\bc(\bx^{k+1})\|
    \le
    C\left(
    \frac{1}{\bar\rho_k}+\|\Delta\bx^{k+1}\|
    \right).
\end{equation}
Using $(u+v)^2\le2u^2+2v^2$ in
\eqref{eq:next-feasibility-from-step}, we obtain
\[
    \|\bc(\bx^{k+1})\|^2
    \le
    C\left(
    \frac{1}{\bar\rho_k^2}+\|\Delta\bx^{k+1}\|^2
    \right).
\]
Since the schedule $\bar\rho_k=\rho+c_s(k+1)^a$ is nondecreasing, we have
$\bar\rho_{k+1}-\bar\rho_k\ge0$.  Hence the penalty drift satisfies
\begin{equation}\label{eq:penalty-drift-decomposition}
\begin{aligned}
    \frac{\bar\rho_{k+1}-\bar\rho_k}{2}
    \|\bc(\bx^{k+1})\|^2
    &\le
    C\frac{\bar\rho_{k+1}-\bar\rho_k}{\bar\rho_k^2}
    +C(\bar\rho_{k+1}-\bar\rho_k)
    \|\Delta\bx^{k+1}\|^2 .
\end{aligned}
\end{equation}

We next record the elementary consequences of the schedule for
$\bar\rho_k$.  Since $t\mapsto t^a$ is concave on $(0,\infty)$ for
$a\in(0,1)$,
\[
    (k+2)^a-(k+1)^a
    \le
    a(k+1)^{a-1}.
\]
Thus
\[
    0\le
    \bar\rho_{k+1}-\bar\rho_k
    \le
    C(k+1)^{a-1}.
\]
Since $\bar\rho_k\ge c_s(k+1)^a$, it follows that
\[
    \frac{\bar\rho_{k+1}-\bar\rho_k}{\bar\rho_k^2}
    \le
    C(k+1)^{-a-1},
    \qquad
    \frac{\bar\rho_{k+1}-\bar\rho_k}{\bar\rho_k}
    \le
    C(k+1)^{-1}.
\]
In particular,
\[
    \sum_{k=0}^{\infty}
    \frac{\bar\rho_{k+1}-\bar\rho_k}{\bar\rho_k^2}
    <
    \infty,
    \qquad
    \frac{\bar\rho_{k+1}-\bar\rho_k}{\bar\rho_k}
    \to0 .
\]
By \eqref{eq:mu-drift-tau} and \eqref{eq:adaptive-anchor-stepsize},
\[
    \|\Delta\bmu^{k+1}\|
    =
    \tau_k\|\bc(\bx^k)\|
    \le
    \frac{B_\tau}{(k+1)^2\log(k+2)},
\]
and therefore
\[
    \sum_{k=0}^{\infty}\|\Delta\bmu^{k+1}\|<\infty.
\]
Combining these estimates gives
\[
\begin{aligned}
    \Phi_{k+1}(\bx^{k+1};\bmu^{k+1})-\Phi_k(\bx^k;\bmu^k)
    &\le
    -\frac{1}{2}\bar\rho_k\|\Delta\bx^{k+1}\|^2
    +C(\bar\rho_{k+1}-\bar\rho_k)
    \|\Delta\bx^{k+1}\|^2\\
    &\quad+
    C\|\Delta\bmu^{k+1}\|
    +C\frac{\bar\rho_{k+1}-\bar\rho_k}{\bar\rho_k^2}.
\end{aligned}
\]
Since $(\bar\rho_{k+1}-\bar\rho_k)/\bar\rho_k\to0$, there exists
$k_0\ge0$ such that, for every $k\ge k_0$,
\[
    C(\bar\rho_{k+1}-\bar\rho_k)
    \le
    \frac{1}{4}\bar\rho_k .
\]
Therefore, for every $k\ge k_0$,
\[
    \Phi_{k+1}(\bx^{k+1};\bmu^{k+1})-\Phi_k(\bx^k;\bmu^k)
    \le
    -\frac{1}{4}\bar\rho_k\|\Delta\bx^{k+1}\|^2
    +C\|\Delta\bmu^{k+1}\|
    +C\frac{\bar\rho_{k+1}-\bar\rho_k}{\bar\rho_k^2}.
\]
For $K>k_0$, summing this inequality from $k=k_0$ to $K-1$ yields
\[
\begin{aligned}
    \frac{1}{4}
    \sum_{k=k_0}^{K-1}
    \bar\rho_k\|\Delta\bx^{k+1}\|^2
    &\le
    \Phi_{k_0}(\bx^{k_0};\bmu^{k_0})
    -
    \Phi_K(\bx^K;\bmu^K)\\
    &\quad+
    C\sum_{k=k_0}^{K-1}\|\Delta\bmu^{k+1}\|
    +
    C\sum_{k=k_0}^{K-1}
    \frac{\bar\rho_{k+1}-\bar\rho_k}{\bar\rho_k^2}.
\end{aligned}
\]
Compactness of $S$, boundedness of $\{\bmu^k\}$, and
$\|\bc(\bx)\|\le C_c$ imply the lower bound
\[
    \begin{aligned}
    \Phi_K(\bx^K;\bmu^K)
    &=
    f(\bx^K)
    +\langle\bmu^K,\bc(\bx^K)\rangle
    +\frac{\bar\rho_K}{2}\|\bc(\bx^K)\|^2\\
    &\ge
    f(\bx^K)-\|\bmu^K\|\,\|\bc(\bx^K)\|
    \ge
    \inf_{\bx\in S}f(\bx)-M_\mu C_c .
    \end{aligned}
\]
Here we used $\frac{\bar\rho_K}{2}\|\bc(\bx^K)\|^2\ge0$.
Using this lower bound and extending the two summable tails to infinity gives
\[
\begin{aligned}
    \frac{1}{4}
    \sum_{k=k_0}^{K-1}
    \bar\rho_k\|\Delta\bx^{k+1}\|^2
    &\le
    \Phi_{k_0}(\bx^{k_0};\bmu^{k_0})
    -
    \left(\inf_{\bx\in S} f(\bx)-M_\mu C_c\right)\\
    &\quad+
    C\sum_{k=k_0}^{\infty}\|\Delta\bmu^{k+1}\|
    +
    C\sum_{k=k_0}^{\infty}
    \frac{\bar\rho_{k+1}-\bar\rho_k}{\bar\rho_k^2}.
\end{aligned}
\]
The finite prefix is controlled by the working-set diameter:
\[
    \sum_{k=0}^{k_0-1}
    \bar\rho_k\|\Delta\bx^{k+1}\|^2
    \le
    D_S^2\sum_{k=0}^{k_0-1}\bar\rho_k
    <\infty .
\]
If $K\le k_0$, the desired bound follows directly from this finite-prefix
estimate.
Thus one may take
\[
\begin{aligned}
    C_{\mathrm{desc}}
    &:=
    D_S^2\sum_{k=0}^{k_0-1}\bar\rho_k\\
    &\quad+
    4\Bigg[
    \Phi_{k_0}(\bx^{k_0};\bmu^{k_0})
    -
    \left(\inf_{\bx\in S} f(\bx)-M_\mu C_c\right)
    +
    C\sum_{k=k_0}^{\infty}\|\Delta\bmu^{k+1}\|\\
    &\qquad\qquad+
    C\sum_{k=k_0}^{\infty}
    \frac{\bar\rho_{k+1}-\bar\rho_k}{\bar\rho_k^2}
    \Bigg].
\end{aligned}
\]
The value $\Phi_{k_0}(\bx^{k_0};\bmu^{k_0})$ is uniformly bounded above by
the compactness of $S$, the bound $\|\bmu^{k_0}\|\le M_\mu$, and
$\|\bc(\bx^{k_0})\|\le C_c$.  Hence $C_{\mathrm{desc}}$ is finite and
independent of $K$.  Combining the tail and finite-prefix bounds proves
\eqref{eq:weighted-finite-length}.
\end{proof}

\begin{remark}
The proof of Lemma~\ref{lem:weighted-finite-length} does not require a
concrete formula for the auxiliary cap $\bar\tau_k$.  It only uses the
finite-variation estimate implied by \eqref{eq:adaptive-anchor-stepsize};
the structural condition $0\le\tau_k\le\bar\tau_k\le s_k$ is kept for the
dual-smoothing interpolation interpretation.
\end{remark}

\begin{theorem}
\label{theorem:target-complexity}
Suppose that Assumptions~\ref{ass:working-set}--\ref{ass:strong-licq} hold
on the working set $S$.  Assume that $\bar\rho_k=\rho+s_k$ satisfies
$\bar\rho_k=\rho+c_s(k+1)^a$ with $c_s>0$ and $a\in(0,1)$.  Let the primal
stepsizes satisfy $\eta_k^{-1}\le C_\eta\bar\rho_k$ and the sufficient
descent condition \eqref{eq:stepsize-sufficient-descent}.  Assume also that
the dual stepsizes satisfy \eqref{eq:adaptive-anchor-stepsize}.

For any $K\ge2$, choose
\[
    \kappa
    \in
    \argmin_{k\in\{\lceil K/2\rceil,\ldots,K-1\}}
    \bar\rho_k\|\Delta\bx^{k+1}\|^2.
\]
Define the effective KKT multiplier
\[
    \by^\kappa
    :=
    \bmu^\kappa+\bar\rho_\kappa\bc(\bx^\kappa).
\]
Then
\begin{equation}\label{eq:kkt-residual-general-a}
    \left\|
    \nabla f(\bx^\kappa)
    +
    J_{\bc}(\bx^\kappa)^\top \by^\kappa
    \right\|
    =
    O\!\left(K^{-(1-a)/2}\right),
    \qquad
    \|\bc(\bx^\kappa)\|
    =
    O(K^{-a}).
\end{equation}
Moreover, $\{\by^\kappa\}$ is bounded.  In particular, for the balanced
choice $a=1/3$,
\[
    \left\|
    \nabla f(\bx^\kappa)
    +
    J_{\bc}(\bx^\kappa)^\top \by^\kappa
    \right\|
    =
    O(K^{-1/3}),
    \qquad
    \|\bc(\bx^\kappa)\|
    =
    O(K^{-1/3}).
\]
Consequently, to obtain an $\epsilon$-approximate KKT point in the sense that
\[
    \left\|
    \nabla f(\bx^\kappa)
    +
    J_{\bc}(\bx^\kappa)^\top \by^\kappa
    \right\|
    \le
    \epsilon,
    \qquad
    \|\bc(\bx^\kappa)\|
    \le
    \epsilon,
\]
under the assumptions above, it suffices to take
\[
    K
    =
    O(\epsilon^{-3}).
\]
Moreover, the auxiliary dual-smoothing multiplier satisfies
\[
    \blmbd^\kappa
    =
    \bmu^\kappa+s_\kappa\bc(\bx^\kappa),
    \qquad
    \by^\kappa
    =
    \blmbd^\kappa+\rho\bc(\bx^\kappa).
\]
Equivalently, $\by^\kappa-\blmbd^\kappa=\rho\bc(\bx^\kappa)$, so their
difference is directly controlled by the feasibility residual.  The KKT
residual is naturally stated with $\by^\kappa$, because it is the multiplier
appearing in the primal stationarity relation.
\end{theorem}

\begin{proof}
Lemma~\ref{lem:weighted-finite-length} gives
\[
    \sum_{k=0}^{K-1}
    \bar\rho_k\|\Delta\bx^{k+1}\|^2
    \le
    C_{\mathrm{desc}},
    \qquad
    \forall K\ge1.
\]
The index set $\{\lceil K/2\rceil,\ldots,K-1\}$ has
$\lfloor K/2\rfloor$ elements, and $\lfloor K/2\rfloor\ge K/3$ for every
$K\ge2$.  Hence, by the definition of $\kappa$,
\[
    \bar\rho_\kappa\|\Delta\bx^{\kappa+1}\|^2
    \le
    \frac{3C_{\mathrm{desc}}}{K}.
\]
Since $\kappa\in\{\lceil K/2\rceil,\ldots,K-1\}$ and
$\bar\rho_k=\rho+c_s(k+1)^a$, we have $\bar\rho_\kappa\asymp K^a$.  Thus
\[
    \|\Delta\bx^{\kappa+1}\|
    =
    O\!\left(K^{-(1+a)/2}\right).
\]
The primal update gives
\[
    \nabla f(\bx^\kappa)
    +
    J_{\bc}(\bx^\kappa)^\top \by^\kappa
    =
    \frac{1}{\eta_\kappa}
    (\bx^\kappa-\bx^{\kappa+1}).
\]
Using $\eta_\kappa^{-1}\le C_\eta\bar\rho_\kappa=O(K^a)$ yields
\[
    \left\|
    \nabla f(\bx^\kappa)
    +
    J_{\bc}(\bx^\kappa)^\top \by^\kappa
    \right\|
    =
    O\!\left(K^{-(1-a)/2}\right).
\]

For feasibility, Lemma~\ref{lem:effective-multiplier-feasibility-bound} gives
\[
    \|\bc(\bx^\kappa)\|
    \le
    \frac{1}{\bar\rho_\kappa}
    \left(
    \frac{M_f}{\nu}+M_\mu
    \right)
    +
    \frac{1}{\eta_\kappa\nu\bar\rho_\kappa}
    \|\Delta\bx^{\kappa+1}\|.
\]
Since $\bar\rho_\kappa\asymp K^a$ and
$\eta_\kappa^{-1}\le C_\eta\bar\rho_\kappa$, this implies
\[
    \|\bc(\bx^\kappa)\|
    =
    O(K^{-a})
    +
    O\!\left(K^{-(1+a)/2}\right)
    =
    O(K^{-a}),
\]
where the last equality uses $a\in(0,1)$.

It remains to record boundedness of the returned multiplier.  By uniform
strong LICQ and the primal stationarity relation,
\[
    \nu\|\by^\kappa\|
    \le
    \|J_{\bc}(\bx^\kappa)^\top\by^\kappa\|
    \le
    M_f
    +
    \left\|
    \nabla f(\bx^\kappa)
    +
    J_{\bc}(\bx^\kappa)^\top\by^\kappa
    \right\|.
\]
The stationarity bound above gives $\|\by^\kappa\|\le C$ uniformly in $K$.
Thus \eqref{eq:kkt-residual-general-a} follows.

The two residual rates are balanced by solving
$a=(1-a)/2$, which gives $a=1/3$. For this choice, both stationarity and
feasibility are $O(K^{-1/3})$, and the iteration complexity claim follows by
setting $K=O(\epsilon^{-3})$.
\end{proof}

\subsection{Dual-damping linearized ALM}

\begin{align*}
    s_k&=c_s(k+1)^{1/3},
    &
    \bar\rho_k&=\rho+s_k,
    &
    \eta_k&=(C_\eta\bar\rho_k)^{-1},\\
    \theta_k
    &=
    \min\left\{
    \frac{B_\tau}
    {s_k\|\bc(\bx^k)\|(k+1)^2\log(k+2)},
    \bar\theta_k
    \right\},
    &
    0&\le\bar\theta_k\le1 .
\end{align*}
where $\rho>0$ is the base penalty, $c_s>0$,
$B_\tau\asymp\max\{1,\|\bc(\bx^0)\|\}$, and $C_\eta$ is chosen large enough
so that \eqref{eq:stepsize-sufficient-descent} holds.
The rate exponent is induced by the growth
$s_k=c_s(k+1)^{1/3}$, while the base penalty $\rho$ affects the constants
through the initial effective penalty $\bar\rho_0=\rho+s_0$ and the
stepsize constant $C_\eta$. 

 
\begin{algorithm}[H]
    \caption{Dual-damping linearized ALM}
    \label{alg:adaptive-dsalm-interpolation}
    \small
    \begin{algorithmic}[1]
        \STATE \textbf{Input:} $\bx^0\in\RR^n$, anchor multiplier
        $\bmu^0\in\RR^m$, base penalty $\rho>0$, smoothing scales
        $\{s_k\}_{k\ge0}$, primal stepsizes $\{\eta_k\}_{k\ge0}$,
        and interpolation weights $\{\theta_k\}_{k\ge0}\subset[0,1]$.
        \STATE Set
        \[
            \blmbd^0=\bmu^0+s_0\bc(\bx^0),
            \qquad
            \by^0=\blmbd^0+\rho\bc(\bx^0)
            =
            \bmu^0+(\rho+s_0)\bc(\bx^0).
        \]
        \FOR{$k=0,1,\ldots,K-1$}
            % \STATE Set the effective multiplier
            % \[
            %     \by^k
            %     =
            %     \blmbd^k+\rho\bc(\bx^k)
            %     =
            %     \bmu^k+\bar\rho_k\bc(\bx^k),
            %     \qquad
            %     \bar\rho_k:=\rho+s_k.
            % \]
            \STATE Update the primal variable
            \[
                \bx^{k+1}
                \in
                \argmin_{\bx\in\RR^n}
                \left\{
                \left\langle
                \nabla f(\bx^k)
                +J_{\bc}(\bx^k)^\top(\blmbd^k+\rho\bc(\bx^k)),
                \bx-\bx^k
                \right\rangle
                +
                \frac{1}{2\eta_k}\|\bx-\bx^k\|^2
                \right\}.
            \]
            \STATE Update the anchor multiplier by interpolation
            \[
                \bmu^{k+1}
                =
                (1-\theta_k)\bmu^k+\theta_k\blmbd^k.
            \]
            \STATE Update the virtual smoothed multiplier
            \[
                \blmbd^{k+1}
                =
                \bmu^{k+1}+s_{k+1}\bc(\bx^{k+1}).
            \]
        \ENDFOR
        \STATE \textbf{Output:} for $K\ge2$, choose
        \[
            \kappa
            \in
            \argmin_{k\in\{\lceil K/2\rceil,\ldots,K-1\}}
            \bar\rho_k\|\bx^{k+1}-\bx^k\|^2,
        \]
        and return $(\bx^\kappa,\bmu^\kappa,\blmbd^\kappa)$.
    \end{algorithmic}
\end{algorithm}
 
\end{document}
